% Lösungen Trigonometrie · Sinussatz · Prüfungs-Fokus MSA (zusammenbau v0.6)
\einheitenkopf{Einheit 4 · Sinussatz}

\erg{1}{a) sin, cos, tan \quad b) $\frac{b}{\mathrm{sin}\,51^\circ} = \frac{7}{\mathrm{sin}\,68^\circ}$, $b = 7 \cdot \mathrm{sin}\,51^\circ : \mathrm{sin}\,68^\circ \approx 5{,}9$ cm \quad c) $PQ$ liegt gegenüber $R$, $PR$ gegenüber $Q$: $PR = 8 \cdot \mathrm{sin}\,53^\circ : \mathrm{sin}\,69^\circ \approx 6{,}8$ cm}
\erg{2}{a) $BD$ liegt gegenüber $118^\circ$, $AD$ gegenüber $33^\circ$: $BD = 4{,}6 \cdot \mathrm{sin}\,118^\circ : \mathrm{sin}\,33^\circ \approx 7{,}5$ cm \quad b) $GS = 3\,200 \cdot \mathrm{sin}\,57^\circ : \mathrm{sin}\,72^\circ \approx 2\,822$ m; Weg $\approx 4\,622$ m $\approx 4{,}6$ km}
\erg{3}{a) $\gamma = 180^\circ - 31^\circ - 113^\circ = 36^\circ$; $BC = 520 \cdot \mathrm{sin}\,31^\circ : \mathrm{sin}\,36^\circ \approx 455{,}6$ m \quad b) $\gamma = 180^\circ - 27^\circ - 121^\circ = 32^\circ$; $BC = 290 \cdot \mathrm{sin}\,27^\circ : \mathrm{sin}\,32^\circ \approx 248{,}4$ m \quad c) dritter Winkel $180^\circ - 6^\circ - 137^\circ = 37^\circ$ (gegenüber $220$ cm); $y = 220 \cdot \mathrm{sin}\,137^\circ : \mathrm{sin}\,37^\circ \approx 249{,}3$ cm \quad d) dritter Winkel $180^\circ - 4^\circ - 148^\circ = 28^\circ$ (gegenüber $190$ cm); $y = 190 \cdot \mathrm{sin}\,148^\circ : \mathrm{sin}\,28^\circ \approx 214{,}5$ cm \quad e) Das Dreieck ist nicht rechtwinklig: Sinussatz. $AC = 9 \cdot \mathrm{sin}\,62^\circ : \mathrm{sin}\,4^\circ \approx 113{,}92$ m, also rund $114$ m \quad f) Das Dreieck ist nicht rechtwinklig: Sinussatz. $AC = 8 \cdot \mathrm{sin}\,51^\circ : \mathrm{sin}\,6^\circ \approx 59{,}48$ m, also rund $59{,}5$ m \quad g) $\angle ACB = 31^\circ$, $\angle ACD = 57^\circ$; $AD = 7 \cdot \mathrm{sin}\,57^\circ : \mathrm{sin}\,78^\circ \approx 6{,}0$ m \quad h) $\angle ACB = 37^\circ$, $\angle ACD = 59^\circ$; $AD = 10 \cdot \mathrm{sin}\,59^\circ : \mathrm{sin}\,69^\circ \approx 9{,}2$ m \quad i) Winkel bei $D$: $39^\circ$; $DE = 9 \cdot \mathrm{sin}\,118^\circ : \mathrm{sin}\,39^\circ \approx 12{,}63$ m; $DP = DE - 8 \approx 4{,}6$ m \quad j) Winkel bei $D$: $46^\circ$; $DE = 16 \cdot \mathrm{sin}\,107^\circ : \mathrm{sin}\,46^\circ \approx 21{,}27$ m; $DP = DE - 17 \approx 4{,}3$ m \quad k) $\angle ABD = 42^\circ$, $\angle DBC = 62^\circ$; $BC = 1\,300 \cdot \mathrm{sin}\,57^\circ : \mathrm{sin}\,62^\circ \approx 1\,235$ m – nicht gleich $1\,300$ m, die Aussage ist falsch. \quad l) $\angle ABD = 39^\circ$, $\angle DBC = 73^\circ$; $BC = 1\,760 \cdot \mathrm{sin}\,63^\circ : \mathrm{sin}\,73^\circ \approx 1\,640$ m – nicht gleich $1\,760$ m, die Aussage ist falsch.}
